\capitulo{CAPÍTULO 10}{La responsabilidad patrimonial del agente estatal: la acción de repetición}

\section{10.1. Normas relevantes anteriores a la Constitución (Código Contencioso Administrativo)}
\begin{cita}
\textit{Art. 77 del CCA: «Sin perjuicio de la responsabilidad que corresponda a la Nación y a las entidades territoriales o descentralizadas, o a las privadas que cumplan funciones públicas, los funcionarios serán responsables de los daños que causen por culpa grave o dolo en el ejercicio de sus funciones.»}
\par
\textit{Art. 78 del CCA: «Los perjudicados podrán demandar, ante la jurisdicción en lo contencioso administrativo según las reglas generales, a la entidad, al funcionario o a ambos. Si prospera la demanda contra la entidad o contra ambos y se considera que el funcionario debe responder, en todo o en parte, la sentencia dispondrá que satisfaga los perjuicios la entidad. En este caso la entidad repetirá contra el funcionario por lo que le correspondiere.»}
\par
\textit{Art. 86 del CCA: «(...) Las entidades públicas deberán promover la misma acción [reparación directa] cuando resulten condenadas o hubieren conciliado por una actuación administrativa originada en culpa grave o dolo de un servidor o ex servidor público que no estuvo vinculado al proceso respectivo (…).»}
\par
\textit{Art. 136-9 del CCA (caducidad): «La de repetición caducará al vencimiento del plazo de dos (2) años, contados a partir del día siguiente de la fecha del pago total efectuado por la entidad.»}
\par
\textit{Sentencia C-832 de 2001: «a partir de la fecha en que efectivamente se realice el pago, o, a más tardar, desde el vencimiento del plazo de 18 meses previsto en el artículo 177 inciso 4 del Código Contencioso Administrativo».}
\end{cita}
\begin{ampliacion}
\textcolor[HTML]{5B3F8C}{\textbf{Ampliación. }}Hoy, la repetición es un medio de control autónomo (art. 142 del CPACA) regulado en lo sustancial y procesal por la Ley 678 de 2001, modificada por la Ley 2195 de 2022. El CCA (Decreto 01 de 1984) ya permitía demandar a la entidad, al funcionario o a ambos (art. 78) y ordenaba a las entidades condenadas repetir (art. 86); la Corte Constitucional, en la C-832 de 2001, condicionó el inicio del término de caducidad de dos años a la fecha del pago efectivo o, a más tardar, al vencimiento del plazo de 18 meses que tenía la entidad para pagar la condena.
\end{ampliacion}
\section{10.2. La Constitución de 1991}
\subsection{10.2.1. La ponencia del artículo 90}
\begin{cita}
\textit{«(…) Finalmente, he estimado de suma utilidad acoger dos planteamientos contenidos en algunos proyectos: el de que, para proteger a las víctimas, se les permita demandar indistintamente al Estado o a la autoridad pública involucrada; y el de que, para proteger el patrimonio público y combatir la indolencia personal del grueso de los funcionarios públicos, se le imponga al Estado el deber de repetir contra el funcionario culpable por las sumas por las que aquél hubiera sido condenado. (…)» (Gaceta Constitucional No. 56, p. 14)}
\end{cita}
Se distinguen así la responsabilidad patrimonial de la entidad estatal y la responsabilidad patrimonial del agente estatal.

\begin{cita}
\textit{«(…) El delegatario Ramírez Ocampo (…) insiste en reducir el artículo, agrega que la repetición contra el funcionario va a conducir a que no haya quien acepte desempeñar cargos públicos, el constituyente Uribe Vargas se muestra de acuerdo. (…)» (Gaceta Constitucional No. 132, p. 19)}
\end{cita}
De allí surge la culpa grave o el dolo como limitación para repetir.

\subsection{10.2.2. El artículo 90 de la Constitución}
\begin{cita}
\textit{«ARTÍCULO 90. El Estado responderá patrimonialmente por los daños antijurídicos que le sean imputables, causados por la acción o la omisión de las autoridades públicas. En el evento de ser condenado el Estado a la reparación patrimonial de uno de tales daños, que haya sido consecuencia de la conducta dolosa o gravemente culposa de un agente suyo, aquél deberá repetir contra éste.»}
\end{cita}
El artículo 90 tiene dos incisos: el primero regula la responsabilidad patrimonial del Estado; el segundo, la responsabilidad patrimonial del agente, es decir, la acción de repetición.

\section{10.3. Limitaciones constitucionales para el ejercicio de la acción de repetición}
Del segundo inciso del artículo 90 se desprenden tres limitaciones (garantías) para exigir la responsabilidad patrimonial del agente: dos procesales y una sustancial.

\subsection{10.3.1. Primera limitación (procesal): la condena previa contra el Estado}
Para que pueda iniciarse la acción de repetición, tiene que haber una condena previa contra el Estado, porque el artículo 90 dice que el Estado repetirá cuando haya sido condenado.

\begin{ejemplo}
\textcolor[HTML]{4F6228}{\textbf{Ejemplo. }}María Victoria, empleada pública del Ministerio de Trabajo, expide actos administrativos (por ejemplo, impone multas por incumplimiento de normas laborales). Tiene una garantía: mientras no se condene al Ministerio de Trabajo, su patrimonio está a salvo; no hay forma de demandarla a ella hasta que exista esa condena. Primero tiene que haber un demandante que demande al Estado y, en ese proceso, una condena. María Victoria no es parte en ese primer proceso: a quien condenan es al Estado.
\end{ejemplo}
\textbf{La condena debe ser indemnizatoria (no restitutoria).} No basta cualquier condena: tiene que ser una condena a reparar un daño antijurídico. Si la condena es restitutoria, la acción de repetición es improcedente.

\begin{ejemplo}
\textcolor[HTML]{4F6228}{\textbf{Ejemplo. }}Caracol, RCN y la ANTV (caso real). Caracol y RCN tenían un contrato de concesión con el Estado y la obligación contractual de transmitir a las siete de la noche un contenido del Ministerio de Cultura; no lo hicieron. La Autoridad Nacional de Televisión (entidad concedente) les impuso a cada uno una multa contractual de 500 millones de pesos, que pagaron. Demandaron ante un tribunal de arbitramiento, que anuló parcialmente la multa: sí había que multarlos, pero el tope permitido era de 200 millones. En consecuencia, el tribunal impuso una condena restitutoria: la ANTV debía devolver 300 millones. La ANTV repitió contra los tres comisionados que firmaron el acto. El Consejo de Estado dijo que la acción era improcedente, porque la condena era restitutoria: si se les cobra a los comisionados, se genera un enriquecimiento para la entidad, que nunca tuvo derecho a esos 300 millones (era un pago en exceso que debía devolver).
\par
\textcolor[HTML]{4F6228}{\textbf{Ejemplo. }}El hecho cumplido en obra pública. Un contratista construyó un kilómetro de carretera que nunca le pagaron; demanda y la entidad es condenada a pagarle 100 millones. La entidad no puede repetir contra el funcionario que no reconoció esa suma, porque de todos modos tenía que pagarla: una condena por desequilibrio no es indemnizatoria, sino restitutoria; es algo que se habría tenido que pagar en todo caso.
\end{ejemplo}
\subsection{10.3.2. Segunda limitación (procesal): la legitimación en la causa por activa}
María Victoria sabe que la única que puede demandarla es la entidad condenada (el Ministerio de Trabajo). El demandante del primer proceso no puede ir contra ella: quien repite es el Estado, por lo que fue condenado.

\subsection{10.3.3. Tercera limitación (sustancial): la culpabilidad del agente}
La acción de repetición se activa por la culpa grave o el dolo del agente estatal. Y aquí es esencial distinguir entre \textbf{culpa} y \textbf{culpabilidad}.

\begin{definicion}
\textcolor[HTML]{1F3864}{\textbf{Definición. }}\textbf{Dolo: }la intención de lesionar a alguien.
\par
\textcolor[HTML]{1F3864}{\textbf{Definición. }}\textbf{Culpa grave: }una negligencia tan grande, tan abrupta, que se asemeja al dolo.
\end{definicion}
\begin{ejemplo}
\textcolor[HTML]{4F6228}{\textbf{Ejemplo. }}María Victoria impone una multa después de un estudio juicioso de la jurisprudencia, pero no sabía que tres días antes la Corte Constitucional había expedido una sentencia que cambió la posición. La persona multada demanda, gana el proceso de nulidad y restablecimiento y condenan al Ministerio. Eso es culpa, no culpabilidad. Igual si había dos posiciones jurisprudenciales (una de la Subsección A y otra de la Subsección B), María Victoria acogió la que más la convencía y el juez acogió la otra: haber escogido una en lugar de la otra es culpa; de ninguna forma es culpabilidad. Es una garantía de la agente: solo se repite por dolo o culpa grave.
\end{ejemplo}
\section{10.4. El esquema constitucional}
\begin{itemize}
\item \textbf{Proceso 1:} la víctima demanda a la entidad causante del daño; condena indemnizatoria (no restitutoria).
\item \textbf{Proceso 2:} la entidad condenada demanda al agente.
\end{itemize}
\textbf{Reglas probatorias:}

\begin{itemize}
\item \textbf{Carga de la prueba:} la entidad condenada tiene la carga de probar la culpabilidad del agente.
\item \textbf{Inoponibilidad:} la sentencia condenatoria y las pruebas practicadas en el primer proceso no son oponibles al agente demandado.
\item \textbf{Juicio de culpabilidad en concreto vs. culpa en abstracto.}
\end{itemize}
\begin{ejemplo}
\textcolor[HTML]{4F6228}{\textbf{Ejemplo. }}El falso positivo grabado (ejemplo inventado). El Ejército es condenado en reparación directa por la ejecución de Pedro en un falso positivo. En la sentencia se dice que la falla del servicio está probada con un video (que obra en el expediente) en el que se ve a los soldados amarrar a Pedro en un callejón y ejecutarlo, para luego presentarlo como guerrillero muerto en combate. El Ejército inicia la acción de repetición contra los tres soldados del video y aporta únicamente la sentencia de reparación directa. ¿Es suficiente? No: como los tres soldados no fueron parte del proceso de reparación directa, nada de lo que allí se hizo les es oponible; no pudieron ejercer defensa ni contradicción. Antes de la Ley 678, las entidades buscaban condenas en repetición con la sola sentencia condenatoria, y la jurisprudencia dijo por mucho tiempo que no era suficiente. ¿Qué debió hacer el Ejército? Pedir el traslado de la prueba (el video) al proceso de repetición, para que los agentes pudieran controvertirla.
\end{ejemplo}
\subsection{10.4.1. ¿Cómo se defiende el agente estatal?}
\begin{itemize}
\item \textbf{Primera defensa: desvirtuar su culpabilidad:} demostrar que no actuó con dolo ni culpa grave.
\item \textbf{Segunda defensa: cuestionar la responsabilidad misma del Estado.} Como la sentencia del primer proceso no le es oponible, el agente puede discutir en la repetición incluso la responsabilidad del Estado.
\end{itemize}
\begin{ejemplo}
\textcolor[HTML]{4F6228}{\textbf{Ejemplo. }}John Jairo, conductor de un vehículo del Estado. Según la sentencia de reparación directa, John Jairo atropelló a la víctima y el Estado fue condenado. En la repetición, John Jairo puede alegar que en realidad hubo un hecho de un tercero (otra persona atropelló a la víctima) o un hecho de la víctima: si en el sector había una cámara de semáforo cuyo video muestra claramente que quien atropelló fue otro civil, y el Estado se defendió mal y no pidió ese video, la mala defensa del Estado no lo vincula. Eso no afecta la indemnización de la víctima, porque esa sentencia hizo tránsito a cosa juzgada; solo afecta la responsabilidad patrimonial del agente.
\end{ejemplo}
\begin{comentario}
\textcolor[HTML]{8A6D3B}{\textbf{Comentario. }}Precisiones. ¿Qué pasa cuando varios funcionarios participaron (quien proyecta, el coordinador, la segunda instancia)? La responsabilidad es subjetiva e individual; el problema usual de las acciones de repetición es que se dirigen solo contra quien firmó. Si se demanda solo al ministro por un acto ilegal, él dirá que no es abogado y que confió en el secretario general, el jurídico o el asesor, y probablemente el Consejo de Estado le dará la razón: hay que ver quién actuó con culpabilidad, es decir, quién, según las circunstancias internas, tenía el conocimiento jurídico, sabía o fue advertido. ¿Qué pasa si la entidad defendió en el proceso que el acto era legal y luego repite contra la funcionaria? No hay problema, porque si la sentencia declaró ilegal el acto, la entidad puede luego repetir. Además, las condenas que dan lugar a repetición pueden provenir de cualquier tipo de proceso (controversias contractuales, nulidad y restablecimiento, reparación directa), no solo de la reparación directa.
\end{comentario}
\section{10.5. La Ley 678 de 2001}
La Ley 678 introdujo dos cambios importantes: el llamamiento en garantía con fines de repetición y las presunciones de culpa grave y dolo.

\subsection{10.5.1. El llamamiento en garantía con fines de repetición}
\begin{itemize}
\item \textbf{Proceso 1:} la víctima demanda a la entidad causante del daño.
\item \textbf{Proceso 1:} la entidad demandada llama en garantía al agente.
\end{itemize}
Además del esquema del artículo 90 (víctima contra el Estado y luego Estado contra el agente), la entidad puede, en el mismo proceso de la víctima contra el Estado, llamar en garantía a María Victoria. En la sentencia se estudia primero la responsabilidad del Estado y luego si está probada la culpabilidad de la agente; si lo está, se dirá que lo que pague el Estado deberá ser reintegrado por ella.

\textbf{Implicaciones procesales:}

\begin{itemize}
\item Quien paga la condena a la víctima sigue siendo la entidad estatal; luego la agente reintegra. El llamamiento no altera ni modifica las relaciones procesales: es una acumulación de pretensiones que da economía procesal, porque se resuelven ambas relaciones en un solo proceso.
\item El agente llamado en garantía puede cuestionar la responsabilidad de la entidad o su responsabilidad personal.
\item \textbf{Limitación en las excepciones que puede formular la entidad demandada (art. 19):} si la entidad llama en garantía al agente, no puede alegar al mismo tiempo la causa extraña. Si alega la causa extraña, dice que no causó el daño; si llama en garantía, reconoce que causó el daño y busca que el agente le reintegre. Las dos posiciones son contradictorias, y formularlas de manera subsidiaria tampoco tiene lógica: se pierde por completo el argumento.
\end{itemize}
\begin{ejemplo}
\textcolor[HTML]{4F6228}{\textbf{Ejemplo. }}En el caso del accidente de tránsito, el municipio de Bucaramanga no puede decir a la vez «fue John Jairo, que iba manejando» y «fue la víctima, que se atravesó». Como dice la expresión popular, no se puede ser «cura y pandillero» al tiempo: hay que escoger.
\end{ejemplo}
\subsection{10.5.2. Las presunciones de culpa grave y dolo}
\begin{definicion}
\textcolor[HTML]{1F3864}{\textbf{Definición. }}\textbf{Presunción: }técnica probatoria con un hecho base y un hecho presumido (hecho indicado), unidos por una inferencia lógica. Si se prueba el hecho base, se tiene por acreditado el hecho presumido.
\end{definicion}
\textbf{Presunciones de dolo (Ley 678 de 2001):}

\begin{itemize}
\item 1. Obrar con desviación de poder.
\item 2. Haber expedido el acto administrativo con vicios en su motivación por inexistencia del supuesto de hecho de la decisión adoptada o de la norma que le sirve de fundamento.
\item 3. Haber expedido el acto administrativo con falsa motivación por desviación de la realidad u ocultamiento de los hechos que sirven de sustento a la decisión de la administración.
\item 4. Haber sido penal o disciplinariamente responsable a título de dolo por los mismos daños que sirvieron de fundamento para la responsabilidad patrimonial del Estado.
\item 5. Haber expedido la resolución, el auto o sentencia manifiestamente contrario a derecho en un proceso judicial.
\end{itemize}
\textbf{Presunciones de culpa grave (Ley 678 de 2001):}

\begin{itemize}
\item 1. Violación manifiesta e inexcusable de las normas de derecho.
\item 2. Carencia o abuso de competencia para proferir la decisión anulada, determinada por error inexcusable.
\item 3. Omisión de las formas sustanciales o de la esencia para la validez de los actos administrativos, determinada por error inexcusable.
\item 4. Violar el debido proceso en lo referente a detenciones arbitrarias y dilación en los términos procesales con detención física o corporal.
\end{itemize}
\textbf{Efecto e implicaciones procesales:}

\begin{itemize}
\item \textbf{Efecto:} inversión de la carga de la prueba. Si se prueba el hecho base, ya no le corresponde a la entidad demostrar la culpa, sino al agente demostrar que no hubo culpa grave o dolo.
\item Estructura de las presunciones: hecho base / hecho indicado / inferencia lógica.
\item La entidad condenada debe invocarlas expresamente en la demanda de repetición.
\item La entidad condenada tiene que probar el hecho base.
\item El agente demandado puede desvirtuar la presunción.
\end{itemize}
\begin{ejemplo}
\textcolor[HTML]{4F6228}{\textbf{Ejemplo. }}María Victoria y la multa a Ricardo. María Victoria le impone a Ricardo una multa de 100 millones por supuestamente incumplir normas laborales. Ricardo demanda la multa alegando falsa motivación y desviación de poder (no es cierto que haya violado esas normas; las pruebas del expediente administrativo lo demuestran; la verdadera razón fue que María Victoria le cogió rabia porque él no le prestó los apuntes de clase). Además de las pretensiones restitutorias, pide perjuicios: su prestigio se vio afectado, sus utilidades disminuyeron y sufrió perjuicios morales. El tribunal anula el acto por falsa motivación y desviación de poder y concede todas las pretensiones restitutorias e indemnizatorias. Luego, el Ministerio de Trabajo repite contra María Victoria con la sentencia de nulidad y restablecimiento como único elemento de juicio. Antes de la Ley 678, esa sentencia no era suficiente, porque no le es oponible. Con la Ley 678 sí lo es, porque demuestra el hecho base de la presunción: el acto se anuló por desviación de poder, y eso permite presumir el dolo. Si María Victoria no presenta pruebas, será condenada; pero puede desvirtuar la presunción (por ejemplo, probando que tenía una buena relación con Ricardo y que nunca hubo falsa motivación ni desviación de poder). Ahora ella tiene que hacer un esfuerzo probatorio, porque se invirtió la carga; antes, si la entidad no invocaba presunciones y solo traía la sentencia, la agente no necesitaba esforzarse.
\end{ejemplo}
\textbf{¿Para qué se crearon las presunciones?} Para facilitarle la vida a la entidad condenada: la mayoría de las veces las entidades perdían los procesos de repetición. La Ley 678 intentó ser «un salvavidas». Aun así, las entidades siguen perdiendo las repeticiones, y son muy pocas las condenas. Hay dos razones:

\begin{itemize}
\item \textbf{Los comités de conciliación no entienden la diferencia entre culpa y culpabilidad.} Demandan en casos como el de María Victoria que escogió una de dos posiciones jurisprudenciales, lo que no tiene sentido. Podría decirse que el 90 \% de las repeticiones se fallan en contra de las entidades porque no prueban la culpabilidad.
\item \textbf{Invocan presunciones cuyo hecho base la sentencia no prueba.} Las entidades invocaban siempre la presunción de «violación manifiesta e inexcusable de las normas de derecho». El juez de la repetición respondía que la sentencia anuló el acto por violación de normas (por ejemplo, por no aplicar la jurisprudencia de una subsección), pero en ninguna parte dijo que esa violación fuera manifiesta e inexcusable: no se probó el hecho base y la entidad perdía.
\end{itemize}
\subsection{10.5.3. Otros aspectos procesales relevantes de la Ley 678 de 2001}
\begin{itemize}
\item \textbf{Art. 8:} legitimación por activa en cabeza de la entidad condenada, el Ministerio Público y el Ministerio de Justicia.
\item \textbf{Art. 11:} caducidad regulada en términos similares a los previstos en el art. 136-9 del CCA; caducidad para el pago en cuotas; C-394 de 2002: caducidad sometida al mismo condicionamiento de la C-832 de 2001.
\item \textbf{Art. 23:} medidas cautelares: embargo y secuestro de bienes del agente demandado; inscripción de la demanda sobre bienes sujetos a registro.
\end{itemize}
\section{10.6. La Ley 2195 de 2022}
La Ley 2195 introdujo muchos cambios en la repetición (ver la relación completa más adelante); dos de ellos merecen especial atención:

\subsection{10.6.1. Eliminación de la limitación procesal para formular el llamamiento en garantía}
Se eliminó la restricción del art. 19 de la Ley 678: ahora la entidad puede llamar en garantía al agente y, simultáneamente, alegar la causa extraña. El legislador pensó que hay que darle la lucha al Estado «con todas las armas, por todas las vías».

\begin{comentario}
\textcolor[HTML]{8A6D3B}{\textbf{Comentario. }}Valoración: es válido pensar así, pero no es coherente. Formular el llamamiento en esos términos es muy débil: «fue culpa de John Jairo, pero, por si las moscas, fue solo culpa de la víctima». Ante esa contestación, el juez concluye que una de las dos defensas está abiertamente equivocada: en lugar de fortalecer la defensa, se debilita.
\end{comentario}
\subsection{10.6.2. Ampliación de las presunciones de culpa grave y dolo}
\textbf{Presunciones de dolo (Ley 2195 de 2022)} —siguen siendo, en lo sustancial, las mismas—:

\begin{itemize}
\item 1. Que el acto administrativo haya sido declarado nulo por desviación de poder, indebida motivación, o falta de motivación, y por falsa motivación.
\item 2. Haber sido penal o disciplinariamente responsable a título de dolo por los mismos daños que sirvieron de fundamento para la responsabilidad patrimonial del Estado.
\item 3. Haber expedido la resolución, el auto o sentencia contrario a derecho en un proceso judicial.
\item 4. Obrar con desviación de poder.
\end{itemize}
\textbf{Presunción de culpa grave (Ley 2195 de 2022):}

\begin{cita}
\textit{«Se presumirá que la conducta del agente del Estado es gravemente culposa cuando el daño es consecuencia de una infracción directa a la Constitución o a la ley o de una inexcusable omisión o extralimitación en el ejercicio de las funciones.»}
\end{cita}
Es una «o» disyuntiva: basta la infracción directa a la Constitución o a la ley.

\begin{comentario}
\textcolor[HTML]{8A6D3B}{\textbf{Comentario. }}Crítica: el hecho base de la presunción de culpa grave es absolutamente amplio, porque el legislador quiere que el Estado gane. Si se anula un acto administrativo, el comité de conciliación dirá inmediatamente que hay que repetir, porque los actos se anulan por violar la ley, es decir, por una infracción directa a la ley: el hecho base queda configurado con cualquier sentencia de nulidad. Cualquier acto que expida María Victoria y sea anulado generará una acción de repetición en su contra. ¿Significa eso que el Estado ganará todas las repeticiones? No: se amplió el hecho base, pero el agente sigue teniendo derecho a desvirtuar la presunción. A María Victoria le costará más defenderse (tendrá que presentar una contestación «de verdad»), pero podrá demostrar, por ejemplo, que en el acto expuso toda la evolución jurisprudencial y explicó por qué prefería una tesis sobre otra; seguramente se concluirá que no se demostró la culpabilidad. La lógica de la Ley 2195 en este punto es cuestionable.
\end{comentario}
\subsection{10.6.3. Otros aspectos procesales relevantes de la Ley 2195 de 2022}
\begin{itemize}
\item \textbf{Art. 8:} legitimación por activa en cabeza de la entidad condenada, el Ministerio Público y la ANDJE.
\item \textbf{Art. 11:} la caducidad se extiende a cinco años.
\item \textbf{Art. 12:} conciliación y reducción de capital según ingresos o patrimonio.
\item \textbf{Art. 13:} acuerdos de pago y reducción de capital.
\item \textbf{Arts. 23 a 25:} se elimina la necesidad de prestar caución; se agrega el embargo de salarios.
\end{itemize}
\begin{ampliacion}
\textcolor[HTML]{5B3F8C}{\textbf{Ampliación. }}Las reformas de la Ley 2195 de 2022 a la Ley 678 de 2001, artículo por artículo:
\end{ampliacion}
\begin{itemize}
\item \textbf{Art. 39 (modifica el art. 5, dolo):} hay dolo cuando el agente actúa con una intención ajena a los fines del servicio público; se presume, entre otros eventos, cuando el acto fue anulado por desviación de poder, indebida, falta o falsa motivación, cuando el agente fue declarado penal o disciplinariamente responsable a título de dolo por los mismos hechos, o cuando expidió una providencia judicial contraria a derecho.
\item \textbf{Art. 40 (modifica el art. 6, culpa grave):} se presume la culpa grave cuando el daño es consecuencia de una infracción directa a la Constitución o a la ley, o de una inexcusable omisión o extralimitación en el ejercicio de las funciones.
\item \textbf{Art. 41 (modifica el art. 8, legitimación):} la entidad debe ejercer la acción dentro de los seis meses siguientes al pago; si no lo hace, pueden ejercerla el Ministerio Público o la Agencia Nacional de Defensa Jurídica del Estado.
\item \textbf{Art. 42 (modifica el art. 11, caducidad):} la acción caduca a los cinco años contados a partir del día siguiente al pago.
\item \textbf{Art. 44 (modifica el art. 19, llamamiento en garantía):} la entidad o el Ministerio Público pueden llamar en garantía al agente identificado, para decidir en el mismo proceso la responsabilidad del Estado y la del agente, en cuaderno separado.
\item \textbf{Arts. 45 a 47 (modifican los arts. 23 a 25, medidas cautelares):} procede el embargo y secuestro de bienes del demandado según el CGP, incluido el embargo de salarios dentro de los límites de la ley laboral, y el juez debe decretar las medidas solicitadas antes de notificar el auto admisorio.
\item \textbf{Art. 48 (modifica el art. 12, conciliación):} permite conciliar fórmulas y plazos de pago, con porcentajes mínimos de pago según los ingresos y el patrimonio del demandado.
\item \textbf{Art. 49 (adiciona el art. 13-1, acuerdos de pago):} ejecutoriada la sentencia, permite acuerdos de pago con condonación parcial del capital y de los intereses según las condiciones económicas del condenado y el plazo de pago.
\end{itemize}
{\small \textcolor[HTML]{6B6B6B}{\textit{Fuente: Ley 2195 de 2022.}}}

\section{10.7. Jurisprudencia relevante reciente sobre la acción de repetición}
\textbf{SU-354 de 2020:}

\begin{itemize}
\item Finalidades preventiva y retributiva de la acción de repetición. Carácter subsidiario, subjetivo y sujeto a criterios de proporcionalidad.
\item «Estricto juicio de atribución de responsabilidad»; derecho de defensa.
\item Configurar el supuesto de hecho de la presunción frente al demandado.
\item Determinar la proporcionalidad de las consecuencias contra el agente.
\end{itemize}
\textbf{SU-259 de 2021:}

\begin{itemize}
\item Juicio de reproche subjetivo contra el agente.
\item Garantías del principio de culpabilidad; superar criterios abstractos como el del «buen padre de familia».
\end{itemize}
\begin{ampliacion}
\textcolor[HTML]{5B3F8C}{\textbf{Ampliación. }}Qué aportan estas sentencias. Ambas son sentencias de unificación de la Corte Constitucional que leen la acción de repetición a la luz del debido proceso (art. 29 C.P.) y del principio de culpabilidad. La \textbf{SU-354 de 2020} destaca que la repetición tiene finalidades preventiva y retributiva, que es subsidiaria (solo procede tras una condena al Estado), subjetiva (exige dolo o culpa grave) y sometida a proporcionalidad; exige un «estricto juicio de atribución de responsabilidad» con plenas garantías de defensa, que el supuesto de hecho de la presunción se configure respecto del agente concreto demandado (no basta que se configure frente a la entidad) y que las consecuencias patrimoniales sean proporcionales. La \textbf{SU-259 de 2021} insiste en que debe hacerse un juicio de reproche subjetivo contra el agente y que no basta el estándar abstracto del «buen padre de familia». Es exactamente la distinción entre culpa y culpabilidad explicada con el caso de María Victoria (sección 10.3.3): la sola ilegalidad del acto (culpa en abstracto) no prueba la culpabilidad del funcionario.
\end{ampliacion}
